% 04-cau-noi.tex -- Bridge boxes for Hochreiter 1991 and Bengio 1994
%
% Scope: half-page each; what they found, why they matter, why neither
% solved the problem (that's chapter 04's job).

\section{Ai đã thấy triệu chứng này trước chúng ta?}
\label{sec:ch02-cau-noi}

Những gì bạn vừa thấy; gradient biến mất theo chiều dài chuỗi, loss không
giảm dù mạng đủ lớn; không phải là phát hiện mới. Hai bài báo đã mô tả nó
từ những năm 1990, trước khi ``deep learning'' thành một từ. Mục này tóm tắt
cả hai. Đây không phải là chương khảo cứu; mỗi bài được một hộp cầu nối, và
lập luận chính của chương nằm ngoài các hộp đó.

\begin{bridge}{Hochreiter 1991}
\index{Hochreiter, Sepp|(}%
\index{vanishing gradient!phát hiện đầu tiên|(}%

Năm 1991, Sepp Hochreiter hoàn thành luận văn diploma tại TU Munich dưới sự
hướng dẫn của Jürgen Schmidhuber. Tựa đề: \enquote{Untersuchungen zu
dynamischen neuronalen Netzen} (Khảo sát về mạng nơ-ron động). Ông huấn luyện
RNN trên các bài toán cần nhớ thông tin qua nhiều bước thời gian và quan sát
thấy một hiện tượng: tín hiệu lỗi từ bước cuối truyền về bước đầu bị suy
giảm theo hàm mũ.

Hochreiter gọi đó là \enquote{vấn đề trễ thời gian dài} (the problem of long
time lags). Ông chưa dùng cụm từ \enquote{vanishing gradient}; cụm từ đó
xuất hiện sau này. Nhưng ông đã mô tả chính xác cơ chế: gradient bị nhân với
đạo hàm của hàm kích hoạt và với ma trận trọng số hồi quy ở mỗi bước lùi,
và nếu các hệ số đó nhỏ hơn 1, tích của chúng tiến về 0.

Luận văn viết bằng tiếng Đức. Cộng đồng nghiên cứu quốc tế không đọc được nó.
Phải sáu năm sau, chính Hochreiter và Schmidhuber mới công bố giải pháp: bộ
nhớ LSTM, với cơ chế \enquote{băng chuyền lỗi hằng số}; chủ đề của chương~4.

\index{Hochreiter, Sepp|)}%
\index{vanishing gradient!phát hiện đầu tiên|)}%
\end{bridge}

\begin{bridge}{Bengio, Simard \& Frasconi 1994}
\index{Bengio, Yoshua|(}%

Ba năm sau Hochreiter, Yoshua Bengio, Patrice Simard và Paolo Frasconi công
bố \enquote{Learning Long-Term Dependencies with Gradient Descent is
Difficult} trên IEEE Transactions on Neural Networks~\autocite{bengio1994}.
Đây là bài báo đầu tiên bằng tiếng Anh mô tả vanishing gradient một cách
toán học chặt chẽ.

Đóng góp chính của bài báo này là hai thứ. Thứ nhất: chứng minh rằng xác
suất để gradient được bảo toàn qua $T$ bước giảm theo hàm mũ của $T$, với
điều kiện các trọng số hồi quy không được ràng buộc đặc biệt. Nói cách khác,
một RNN được khởi tạo ngẫu nhiên \emph{gần như chắc chắn} sẽ bị vanishing
gradient nếu chuỗi đủ dài. Thứ hai: xác nhận bằng thực nghiệm trên bài toán
parity và adding problem.

Bài báo kết luận rằng gradient descent \emph{không đủ} để huấn luyện RNN học
phụ thuộc dài hạn. Cần thay đổi kiến trúc hoặc thuật toán huấn luyện. Lời
kêu gọi đó được trả lời ba năm sau; cũng bởi Hochreiter và Schmidhuber, với
LSTM.

\index{Bengio, Yoshua|)}%
\end{bridge}

Cả hai bài báo đều chẩn đoán đúng bệnh, nhưng không ai trong số họ chữa được
nó. Hochreiter biết gradient biến mất nhưng chưa có giải pháp. Bengio chứng
minh điều đó là tất yếu về mặt xác suất, nhưng cũng dừng ở chẩn đoán. Phải
đến chương~4, với sự ra đời của LSTM, vấn đề này mới có một câu trả lời bằng
kiến trúc.

Còn trước mắt, chương tiếp theo; chương~3; sẽ cho bạn công cụ toán học để
\emph{tự} chẩn đoán vanishing và exploding gradient, dựa trên bài báo của
Pascanu, Mikolov và Bengio~2013.
